\chapter{Anterior Cruciate Ligament (ACL) Reconstruction}
\section*{Overview}
ACL reconstruction replaces a torn anterior cruciate ligament with a tendon graft to restore knee stability.
\section*{Why It Is Done}
It may be considered for active patients with instability, combined knee injuries, or activity goals that cannot be met with rehabilitation alone.
\section*{How It Is Performed}
Usually arthroscopically, the surgeon harvests or uses a donor tendon, drills bone tunnels, and secures the graft in the femur and tibia.
\section*{Risks and Recovery}
Infection, stiffness, blood clots, graft failure, nerve injury, and recurrent instability are possible. Rehabilitation is prolonged and return to sport is gradual.
\section*{References}
\begin{enumerate}\item MedlinePlus. ACL Reconstruction. U.S. National Library of Medicine; 2025.\item Current Medical Diagnosis and Treatment. McGraw-Hill; 2025.\end{enumerate}
\clearpage
